\begin{helper}
Let \(C\) be a history cell and let \(T\) be an event.  For each sample, let
\(u_R(\omega)\) and \(u_B(\omega)\) be the numbers of present red and blue
staged-open terminal coordinates of the \(k\)-set \(I\).  Assume
\(|I|=k\), \(0\le p\le1\), and that for every ordered count pair
\((u_R,u_B)\),
\[
\Pr(T\cap\{u_R(\omega)=u_R,\ u_B(\omega)=u_B\}\mid C)\le M
\]
with \(M\ge0\).  Then
\[
\Pr(T\mid C)\le \max(1,k^4)M.
\]
\end{helper}
\begin{proof}
The count events
\[
\{u_R(\omega)=r,\ u_B(\omega)=b\}
\]
are pairwise disjoint as \((r,b)\) varies, and their union contains every
sample in \(T\).  Conditional probability is computed from nonnegative sample
weights, so finite subadditivity gives an upper bound by the sum of the
conditional probabilities over all possible count pairs.

It remains to count the possible pairs.  Each present red staged-open
terminal coordinate is one projected unordered pair from the \(k\)-element set
\(I\), and similarly for blue.  Therefore each of \(u_R(\omega)\) and
\(u_B(\omega)\) has at most \(\binom{k}{2}+1\) possible values.  If \(k=0\),
there is only the count pair \((0,0)\), and the factor \(\max(1,k^4)\) is
\(1\).  If \(k>0\), then \(\binom{k}{2}+1\le k^2\): it is immediate for
\(k=1\), and for \(k\ge2\) follows from
\(\binom{k}{2}+1\le \binom{k}{2}+k\le k^2\).  Hence there are at most
\(k^4\) ordered count pairs when \(k>0\), and at most \(\max(1,k^4)\) in all
cases.  Summing the uniform bound \(M\) over these possible pairs proves the
claim.
\end{proof}
